% Bagian: second-order-logic
% Bab: metatheory
% Subbagian: second-order-arithmetic

\documentclass[../../../include/open-logic-section]{subfiles}

\begin{document}

\olfileid[id]{sol}{met}{spa}

\olsection{Aritmetika Orde Kedua}

Ingat bahwa teori $\Th{PA}$ dari aritmetika Peano mencakup delapan aksioma
dari~$\Th{Q}$,
\begin{align*}
& \lforall[x][\eq/[x'][\Obj 0]]\\
& \lforall[x][\lforall[y][(\eq[x'][y'] \lif \eq[x][y])]]\\
& \lforall[x][(\eq[x][\Obj 0] \lor \lexists[y][\eq[x][y']])]\\ 
& \lforall[x][\eq[(x + \Obj 0)][x]]\\
& \lforall[x][\lforall[y][\eq[(x + y')][(x + y)']]]\\
& \lforall[x][\eq[(x \times \Obj 0)][\Obj 0]]\\
& \lforall[x][\lforall[y][\eq[(x \times y')][((x \times y) + x)]]]\\
& \lforall[x][\lforall[y][(x < y \liff \lexists[z][\eq[(z' + x)][y]])]]\\
\intertext{bersama semua kalimat berbentuk}
& (!A(\Obj 0) \land \lforall[x][(!A(x) \lif !A(x'))]) \lif \lforall[x][!A(x)].
\end{align*}
Yang terakhir merupakan suatu ``skema'', yakni pola yang menghasilkan tak
hingga banyak !!{sentence} dalam bahasa aritmetika, satu untuk setiap
!!{formula}~$!A(x)$. Skema ini kita sebut \emph{skema aksioma induksi}
(orde pertama).
Dalam aritmetika Peano \emph{orde kedua}~$\Th{PA^2}$, induksi dapat dinyatakan
sebagai satu kalimat. $\Th{PA^2}$ terdiri atas delapan aksioma pertama di atas
bersama \emph{aksioma induksi} (orde kedua):
\[
\lforall[X][((X(\Obj 0) \land \lforall[x][(X(x) \lif X(x'))]) \lif \lforall[x][X(x)])].
\]
Aksioma tersebut menyatakan bahwa jika suatu himpunan bagian~$X$ dari
!!{domain} memuat~$\Assign{\Obj{0}}{M}$ dan, bersama setiap~$x \in
\Domain{M}$, juga memuat~$\Assign{\prime}{M}(x)$ (yakni, himpunan itu
``tertutup terhadap penerus''), maka himpunan itu memuat segala sesuatu dalam
!!{domain} (yakni, $X = \Domain{M})$.

Aksioma induksi menjamin bahwa setiap !!{structure} yang memenuhinya hanya
memuat !!{element}s dari~$\Domain{M}$ yang keberadaannya dituntut oleh
aksioma-aksioma tersebut, yakni nilai-nilai~$\num{n}$ untuk $n \in \Nat$.
Suatu model dari $\Th{PA^2}$ tidak memuat bilangan nonstandar.

\begin{thm}
\ollabel{thm:sol-pa-standard}
Jika $\Sat{M}{\Th{PA^2}}$, maka $\Domain{M} =
\Setabs{\Value{\num{n}}{M}}{n \in \Nat}$.
\end{thm}

\begin{proof}
Misalkan $N = \Setabs{\Value{\num{n}}{M}}{n \in \Nat}$, dan andaikan
$\Sat{M}{\Th{PA^2}}$. Tentu saja, untuk setiap $n \in \Nat$,
$\Value{\num{n}}{M} \in \Domain{M}$, sehingga $N \subseteq \Domain{M}$.

Sekarang tinjau inklusi ke arah sebaliknya. Tinjau suatu penugasan variabel
$s$ dengan $s(X) = N$. Berdasarkan asumsi,
\begin{align*}
\Sat{M}{\lforall[X][((X(\Obj 0) \land \lforall[x][(X(x)
      \lif X(x'))]) \lif \lforall[x][X(x)])]}, & \text{ sehingga}\\
\Sat{M}{(X(\Obj 0) \land \lforall[x][(X(x) \lif X(x'))]) \lif
  \lforall[x][X(x)]}[s]. & 
\end{align*}
Tinjau anteseden implikasi ini. $\Value{\Obj{0}}{M} \in N$, sehingga
$\Sat{M}{X(\Obj{0})}[s]$. Konjungsi kedua,
$\lforall[x][(X(x) \lif X(x'))]$, juga dipenuhi. Sebab, andaikan $x \in N$.
Menurut definisi~$N$, $x = \Value{\num{n}}{M}$ untuk suatu~$n$. Hal ini
memberikan $\Assign{\prime}{M}(x) = \Value{\num{n+1}}{M} \in N$. Jadi,
$\Assign{\prime}{M}(x) \in N$.

Kita memperoleh $\Sat{M}{X(\Obj 0) \land \lforall[x][(X(x) \lif X(x'))]}[s]$.
Akibatnya, $\Sat{M}{\lforall[x][X(x)]}[s]$. Namun, ini berarti bahwa untuk
setiap $x \in \Domain{M}$ berlaku $x \in s(X) = N$. Jadi, $\Domain{M}
\subseteq N$.
\end{proof}

\begin{cor}
\ollabel{cor:sol-pa-categorical}
Sembarang dua model dari~$\Th{PA^2}$ saling isomorfik.
\end{cor}

\begin{proof}
Menurut \olref{thm:sol-pa-standard}, domain setiap model dari $\Th{PA^2}$
seluruhnya terdiri atas $\Value{\num{n}}{M}$. Setiap model semacam itu juga
merupakan model dari~$\Th{Q}$. Menurut
\olref[mod][mar][stm]{prop:thq-standard}, setiap model semacam itu standar,
yakni isomorfik dengan~$\Struct{N}$.
\end{proof}

Di atas kita mendefinisikan $\Th{PA^2}$ sebagai teori yang memuat delapan
aksioma aritmetika pertama bersama aksioma induksi orde kedua. Sebenarnya,
berkat daya ekspresif logika orde kedua, hanya \emph{dua aksioma pertama}
dari aksioma-aksioma aritmetika itu bersama induksi yang diperlukan bagi
aritmetika Peano orde kedua.

\begin{prop}\ollabel{prop:sol-pa-definable}
Misalkan $\Th{PA^{2\dagger}}$ adalah teori orde kedua yang memuat dua
aksioma aritmetika pertama (aksioma-aksioma penerus) dan aksioma induksi
orde kedua. Maka $\le$, $+$, dan $\times$ dapat didefinisikan dalam
$\Th{PA^{2\dagger}}$.
\end{prop}

\begin{proof}
Untuk menunjukkan bahwa $\le$ dapat didefinisikan, kita harus menemukan
suatu formula~$!A_\le(x, y)$ sedemikian sehingga
$\Sat{N}{!A_\le(\num n, \num m)}$ jika dan hanya jika $n \le m$. Tinjau
formula
\[
!B(x, Y) \ident Y(x) \land \lforall[y][(Y(y) \lif Y(y'))]
\]
Jelas bahwa $!B(\num n, Y)$ dipenuhi oleh suatu himpunan $Y \subseteq \Nat$
jika dan hanya jika $\Setabs{m}{n \le m} \subseteq Y$, sehingga kita dapat
mengambil $!A_\le(x, y) \ident \lforall[Y][(!B(x, Y) \lif Y(y))]$.

Untuk melihat bahwa penjumlahan dapat didefinisikan, amati bahwa $k+l=m$
jika dan hanya jika terdapat suatu fungsi $u$ sedemikian sehingga $u(0)=k$,
$u(n')=u(n)'$ untuk semua $n$, dan $m = u(l)$. Kita dapat menggunakan
ekuivalensi ini untuk mendefinisikan penjumlahan dalam
$\Th{PA^{2\dagger}}$ dengan formula berikut:
\[
!A_+(x,y,z) \ident \lexists[u][(u(\Obj 0)=x \land \lforall[w][u(w')=u(w)']
 \land u(y)=z)]
\] 
Harus jelas bahwa $\Sat{N}{!A_+(\num k, \num l, \num m)}$ jika dan hanya
jika $k+l=m$.
\end{proof}

\begin{prob}
Lengkapi pembuktian \olref[sol][met][spa]{prop:sol-pa-definable}.
\end{prob}

\end{document}
